\Sec{General Functions}{Math.GeneralFunctions}

\p This subclause lists and describes functions for general mathematics
operations operating objects of integer types, floating-point types, and vectors
of such types.

\begin{HLSL}
float16_t abs(float16_t x);
float abs(float x);
double abs(double x);
int16_t abs(int16_t x);
uint16_t abs(uint16_t x);
int abs(int x);
uint abs(uint x);
int64_t abs(int64_t x);
uint64_t abs(uint64_t x);
template <uint N> vector<float16_t, N> abs(vector<float16_t, N> x);
template <uint N> vector<float, N> abs(vector<float, N> x);
template <uint N> vector<double, N> abs(vector<double, N> x);
template <uint N> vector<int16_t, N> abs(vector<int16_t, N> x);
template <uint N> vector<uint16_t, N> abs(vector<uint16_t, N> x);
template <uint N> vector<int, N> abs(vector<int, N> x);
template <uint N> vector<uint, N> abs(vector<uint, N> x);
template <uint N> vector<int64_t, N> abs(vector<int64_t, N> x);
template <uint N> vector<uint64_t, N> abs(vector<uint64_t, N> x);

float16_t ceil(float16_t x);
float ceil(float x);
template <uint N> vector<float16_t, N> ceil(vector<float16_t, N> x);
template <uint N> vector<float, N> ceil(vector<float, N> x);

float16_t clamp(float16_t x, float16_t min, float16_t max);
float clamp(float x, float min, float max);
double clamp(double x, double min, double max);
int16_t clamp(int16_t x, int16_t min, int16_t max);
uint16_t clamp(uint16_t x, uint16_t min, uint16_t max);
int clamp(int x, int min, int max);
uint clamp(uint x, uint min, uint max);
int64_t clamp(int64_t x, int64_t min, int64_t max);
uint64_t clamp(uint64_t x, uint64_t min, uint64_t max);
template <uint N>
vector<float16_t, N> clamp(vector<float16_t, N> x, vector<float16_t, N> min,
                           vector<float16_t, N> max);
template <uint N>
vector<float, N> clamp(vector<float, N> x, vector<float, N> min,
                       vector<float, N> max);
template <uint N>
vector<double, N> clamp(vector<double, N> x, vector<double, N> min,
                        vector<double, N> max);
template <uint N>
vector<int16_t, N> clamp(vector<int16_t, N> x, vector<int16_t, N> min,
                         vector<int16_t, N> max);
template <uint N>
vector<uint16_t, N> clamp(vector<uint16_t, N> x, vector<uint16_t, N> min,
                          vector<uint16_t, N> max);
template <uint N>
vector<int, N> clamp(vector<int, N> x, vector<int, N> min, vector<int, N> max);
template <uint N>
vector<uint, N> clamp(vector<uint, N> x, vector<uint, N> min,
                      vector<uint, N> max);
template <uint N>
vector<int64_t, N> clamp(vector<int64_t, N> x, vector<int64_t, N> min,
                         vector<int64_t, N> max);
template <uint N>
vector<uint64_t, N> clamp(vector<uint64_t, N> x, vector<uint64_t, N> min,
                          vector<uint64_t, N> max);

float16_t exp(float16_t x);
float exp(float x);
template <uint N> vector<float16_t, N> exp(vector<float16_t, N> x);
template <uint N> vector<float, N> exp(vector<float, N> x);

float16_t exp2(float16_t x);
float exp2(float x);
template <uint N> vector<float16_t, N> exp2(vector<float16_t, N> x);
template <uint N> vector<float, N> exp2(vector<float, N> x);

float16_t floor(float16_t x);
float floor(float x);
template <uint N> vector<float16_t, N> floor(vector<float16_t, N> x);
template <uint N> vector<float, N> floor(vector<float, N> x);

double fma(double a, double b, double c);
template <uint N>
vector<double, N> fma(vector<double, N> a, vector<double, N> b,
                      vector<double, N> c);

float16_t fmod(float16_t a, float16_t b);
float fmod(float a, float b);
template <uint N>
vector<float16_t, N> fmod(vector<float16_t, N> a, vector<float16_t, N> b);
template <uint N> vector<float, N> fmod(vector<float, N> a, vector<float, N> b);

float16_t frac(float16_t x);
float frac(float x);
template <uint N> vector<float16_t, N> frac(vector<float16_t, N> x);
template <uint N> vector<float, N> frac(vector<float, N> x);

float frexp(float x, out float exp);
template <uint N>
vector<float, N> frexp(vector<float, N> x, out vector<float, N> exp);

float16_t ldexp(float16_t x, float16_t exp);
float ldexp(float x, float exp);
template <uint N>
vector<float16_t, N> ldexp(vector<float16_t, N> x, vector<float16_t, N> exp);
template <uint N>
vector<float, N> ldexp(vector<float, N> x, vector<float, N> exp);

float16_t lerp(float16_t a, float16_t b, float16_t s);
float lerp(float a, float b, float s);
template <uint N>
vector<float16_t, N> lerp(vector<float16_t, N> a, vector<float16_t, N> b,
                          vector<float16_t, N> s);
template <uint N>
vector<float, N> lerp(vector<float, N> a, vector<float, N> b,
                      vector<float, N> s);

float16_t log(float16_t x);
float log(float x);
template <uint N> vector<float16_t, N> log(vector<float16_t, N> x);
template <uint N> vector<float, N> log(vector<float, N> x);

float16_t log10(float16_t x);
float log10(float x);
template <uint N> vector<float16_t, N> log10(vector<float16_t, N> x);
template <uint N> vector<float, N> log10(vector<float, N> x);

float16_t log2(float16_t x);
float log2(float x);
template <uint N> vector<float16_t, N> log2(vector<float16_t, N> x);
template <uint N> vector<float, N> log2(vector<float, N> x);

float16_t mad(float16_t a, float16_t b, float16_t c);
float mad(float a, float b, float c);
double mad(double a, double b, double c);
int16_t mad(int16_t a, int16_t b, int16_t c);
uint16_t mad(uint16_t a, uint16_t b, uint16_t c);
int mad(int a, int b, int c);
uint mad(uint a, uint b, uint c);
int64_t mad(int64_t a, int64_t b, int64_t c);
uint64_t mad(uint64_t a, uint64_t b, uint64_t c);
template <uint N>
vector<float16_t, N> mad(vector<float16_t, N> a, vector<float16_t, N> b,
                         vector<float16_t, N> c);
template <uint N>
vector<float, N> mad(vector<float, N> a, vector<float, N> b,
                     vector<float, N> c);
template <uint N>
vector<double, N> mad(vector<double, N> a, vector<double, N> b,
                      vector<double, N> c);
template <uint N>
vector<int16_t, N> mad(vector<int16_t, N> a, vector<int16_t, N> b,
                       vector<int16_t, N> c);
template <uint N>
vector<uint16_t, N> mad(vector<uint16_t, N> a, vector<uint16_t, N> b,
                        vector<uint16_t, N> c);
template <uint N>
vector<int, N> mad(vector<int, N> a, vector<int, N> b, vector<int, N> c);
template <uint N>
vector<uint, N> mad(vector<uint, N> a, vector<uint, N> b, vector<uint, N> c);
template <uint N>
vector<int64_t, N> mad(vector<int64_t, N> a, vector<int64_t, N> b,
                       vector<int64_t, N> c);
template <uint N>
vector<uint64_t, N> mad(vector<uint64_t, N> a, vector<uint64_t, N> b,
                        vector<uint64_t, N> c);

float16_t max(float16_t a, float16_t b);
float max(float a, float b);
double max(double a, double b);
int16_t max(int16_t a, int16_t b);
uint16_t max(uint16_t a, uint16_t b);
int max(int a, int b);
uint max(uint a, uint b);
int64_t max(int64_t a, int64_t b);
uint64_t max(uint64_t a, uint64_t b);
template <uint N>
vector<float16_t, N> max(vector<float16_t, N> a, vector<float16_t, N> b);
template <uint N> vector<float, N> max(vector<float, N> a, vector<float, N> b);
template <uint N>
vector<double, N> max(vector<double, N> a, vector<double, N> b);
template <uint N>
vector<int16_t, N> max(vector<int16_t, N> a, vector<int16_t, N> b);
template <uint N>
vector<uint16_t, N> max(vector<uint16_t, N> a, vector<uint16_t, N> b);
template <uint N> vector<int, N> max(vector<int, N> a, vector<int, N> b);
template <uint N> vector<uint, N> max(vector<uint, N> a, vector<uint, N> b);
template <uint N>
vector<int64_t, N> max(vector<int64_t, N> a, vector<int64_t, N> b);
template <uint N>
vector<uint64_t, N> max(vector<uint64_t, N> a, vector<uint64_t, N> b);

float16_t min(float16_t a, float16_t b);
float min(float a, float b);
double min(double a, double b);
int16_t min(int16_t a, int16_t b);
uint16_t min(uint16_t a, uint16_t b);
int min(int a, int b);
uint min(uint a, uint b);
int64_t min(int64_t a, int64_t b);
uint64_t min(uint64_t a, uint64_t b);
template <uint N>
vector<float16_t, N> min(vector<float16_t, N> a, vector<float16_t, N> b);
template <uint N> vector<float, N> min(vector<float, N> a, vector<float, N> b);
template <uint N>
vector<double, N> min(vector<double, N> a, vector<double, N> b);
template <uint N>
vector<int16_t, N> min(vector<int16_t, N> a, vector<int16_t, N> b);
template <uint N>
vector<uint16_t, N> min(vector<uint16_t, N> a, vector<uint16_t, N> b);
template <uint N> vector<int, N> min(vector<int, N> a, vector<int, N> b);
template <uint N> vector<uint, N> min(vector<uint, N> a, vector<uint, N> b);
template <uint N>
vector<int64_t, N> min(vector<int64_t, N> a, vector<int64_t, N> b);
template <uint N>
vector<uint64_t, N> min(vector<uint64_t, N> a, vector<uint64_t, N> b);

float16_t pow(float16_t x, float16_t y);
float pow(float x, float y);
template <uint N>
vector<float16_t, N> pow(vector<float16_t, N> x, vector<float16_t, N> y);
template <uint N> vector<float, N> pow(vector<float, N> x, vector<float, N> y);

float16_t rcp(float16_t x);
float rcp(float x);
double rcp(double x);
template <uint N> vector<float16_t, N> rcp(vector<float16_t, N> x);
template <uint N> vector<float, N> rcp(vector<float, N> x);
template <uint N> vector<double, N> rcp(vector<double, N> x);

float16_t round(float16_t x);
float round(float x);
template <uint N> vector<float16_t, N> round(vector<float16_t, N> x);
template <uint N> vector<float, N> round(vector<float, N> x);

float16_t rsqrt(float16_t x);
float rsqrt(float x);
template <uint N> vector<float16_t, N> rsqrt(vector<float16_t, N> x);
template <uint N> vector<float, N> rsqrt(vector<float, N> x);

int sign(float16_t x);
int sign(float x);
int sign(double x);
int sign(int16_t x);
int sign(uint16_t x);
int sign(int x);
int sign(uint x);
int sign(int64_t x);
int sign(uint64_t x);
template <uint N> vector<int, N> sign(vector<float16_t, N> x);
template <uint N> vector<int, N> sign(vector<float, N> x);
template <uint N> vector<int, N> sign(vector<double, N> x);
template <uint N> vector<int, N> sign(vector<int16_t, N> x);
template <uint N> vector<int, N> sign(vector<uint16_t, N> x);
template <uint N> vector<int, N> sign(vector<int, N> x);
template <uint N> vector<int, N> sign(vector<uint, N> x);
template <uint N> vector<int, N> sign(vector<int64_t, N> x);
template <uint N> vector<int, N> sign(vector<uint64_t, N> x);

float16_t smoothstep(float16_t a, float16_t b, float16_t x);
float smoothstep(float a, float b, float x);
template <uint N>
vector<float16_t, N> smoothstep(vector<float16_t, N> a, vector<float16_t, N> b,
                                vector<float16_t, N> x);
template <uint N>
vector<float, N> smoothstep(vector<float, N> a, vector<float, N> b,
                            vector<float, N> x);

float16_t sqrt(float16_t x);
float sqrt(float x);
template <uint N> vector<float16_t, N> sqrt(vector<float16_t, N> x);
template <uint N> vector<float, N> sqrt(vector<float, N> x);

float16_t step(float16_t a, float16_t x);
float step(float a, float x);
template <uint N>
vector<float16_t, N> step(vector<float16_t, N> a, vector<float16_t, N> x);
template <uint N> vector<float, N> step(vector<float, N> a, vector<float, N> x);

float16_t trunc(float16_t x);
float trunc(float x);
template <uint N> vector<float16_t, N> trunc(vector<float16_t, N> x);
template <uint N> vector<float, N> trunc(vector<float, N> x);
\end{HLSL}

\Sub{Other Non-Member Functions}{Math.GeneralFunctions.Other}

\begin{HLSL}
float16_t abs(float16_t x);
float abs(float x);
double abs(double x);
int16_t abs(int16_t x);
uint16_t abs(uint16_t x);
int abs(int x);
uint abs(uint x);
int64_t abs(int64_t x);
uint64_t abs(uint64_t x);
template <uint N> vector<float16_t, N> abs(vector<float16_t, N> x);
template <uint N> vector<float, N> abs(vector<float, N> x);
template <uint N> vector<double, N> abs(vector<double, N> x);
template <uint N> vector<int16_t, N> abs(vector<int16_t, N> x);
template <uint N> vector<uint16_t, N> abs(vector<uint16_t, N> x);
template <uint N> vector<int, N> abs(vector<int, N> x);
template <uint N> vector<uint, N> abs(vector<uint, N> x);
template <uint N> vector<int64_t, N> abs(vector<int64_t, N> x);
template <uint N> vector<uint64_t, N> abs(vector<uint64_t, N> x);
\end{HLSL}

\begin{itemize}
  \item \textit{Returns:} A scalar or vector object of the same dimensions as the
  input where the value of each returned element is the magnitude of the
  corresponding input element (\(|x_0|, .. |x_n|\)).
\end{itemize}

\begin{HLSL}
float16_t ceil(float16_t x);
float ceil(float x);
template <uint N> vector<float16_t, N> ceil(vector<float16_t, N> x);
template <uint N> vector<float, N> ceil(vector<float, N> x);
\end{HLSL}

\begin{itemize}
  \item \textit{Returns:} A scalar or vector object of the same dimensions as the
  input where the value of each returned element is the ceiling of the
  corresponding input element (\(\lceil x_0 \rceil, .. \lceil x_n \rceil\)).
\end{itemize}

\begin{HLSL}
float16_t clamp(float16_t x, float16_t min, float16_t max);
float clamp(float x, float min, float max);
double clamp(double x, double min, double max);
int16_t clamp(int16_t x, int16_t min, int16_t max);
uint16_t clamp(uint16_t x, uint16_t min, uint16_t max);
int clamp(int x, int min, int max);
uint clamp(uint x, uint min, uint max);
int64_t clamp(int64_t x, int64_t min, int64_t max);
uint64_t clamp(uint64_t x, uint64_t min, uint64_t max);
template <uint N>
vector<float16_t, N> clamp(vector<float16_t, N> x, vector<float16_t, N> min,
                           vector<float16_t, N> max);
template <uint N>
vector<float, N> clamp(vector<float, N> x, vector<float, N> min,
                       vector<float, N> max);
template <uint N>
vector<double, N> clamp(vector<double, N> x, vector<double, N> min,
                        vector<double, N> max);
template <uint N>
vector<int16_t, N> clamp(vector<int16_t, N> x, vector<int16_t, N> min,
                         vector<int16_t, N> max);
template <uint N>
vector<uint16_t, N> clamp(vector<uint16_t, N> x, vector<uint16_t, N> min,
                          vector<uint16_t, N> max);
template <uint N>
vector<int, N> clamp(vector<int, N> x, vector<int, N> min, vector<int, N> max);
template <uint N>
vector<uint, N> clamp(vector<uint, N> x, vector<uint, N> min,
                      vector<uint, N> max);
template <uint N>
vector<int64_t, N> clamp(vector<int64_t, N> x, vector<int64_t, N> min,
                         vector<int64_t, N> max);
template <uint N>
vector<uint64_t, N> clamp(vector<uint64_t, N> x, vector<uint64_t, N> min,
                          vector<uint64_t, N> max);
\end{HLSL}

\begin{itemize}
  \item \textit{Returns:} A scalar or vector object of the same dimensions as
  the input where the value of each returned element is clamped between the
  corresponding value in the \texttt{min} and \texttt{max} inputs.
  \item \textit{Remarks:} Behavior of NaN input values is undefined.
\end{itemize}

\begin{HLSL}
float16_t exp(float16_t x);
float exp(float x);
template <uint N> vector<float16_t, N> exp(vector<float16_t, N> x);
template <uint N> vector<float, N> exp(vector<float, N> x);
\end{HLSL}

\begin{itemize}
  \item \textit{Returns:} A scalar or vector object of the same dimensions as
  the input where the value of each returned element is the base $e$ exponent of
  the corresponding element (\(e^{x_0}, .. e^{x_n}\)).
\end{itemize}

\begin{HLSL}
float16_t exp2(float16_t x);
float exp2(float x);
template <uint N> vector<float16_t, N> exp2(vector<float16_t, N> x);
template <uint N> vector<float, N> exp2(vector<float, N> x);
\end{HLSL}

\begin{itemize}
  \item \textit{Returns:} A scalar or vector object of the same dimensions as
  the input where the value of each returned element is the base 2 exponent of
  the corresponding element (\(2^{x_0}, .. 2^{x_n}\)).
\end{itemize}

\begin{HLSL}
float16_t floor(float16_t x);
float floor(float x);
template <uint N> vector<float16_t, N> floor(vector<float16_t, N> x);
template <uint N> vector<float, N> floor(vector<float, N> x);
\end{HLSL}

\begin{itemize}
  \item \textit{Returns:} A scalar or vector object of the same dimensions as
  the input where the value of each returned element is the floor of the
  corresponding input element (\(\lfloor x_0 \rfloor, .. \lfloor x_n \rfloor\)).
\end{itemize}

\begin{HLSL}
double fma(double a, double b, double c);
template <uint N>
vector<double, N> fma(vector<double, N> a, vector<double, N> b,
                      vector<double, N> c);
\end{HLSL}

\begin{itemize}
  \item \textit{Returns:} A scalar or vector object of the same dimensions as
  the input where the value of each returned element is a fused multiply add of
  the corresponding input elements (\(a_0 * b_0 + c_0, .. a_n * b_n + c_n\)).
  \item \textit{Remarks:} The returned value must be accurate to 0.5 units of
  least precision (ULP).
\end{itemize}

\begin{HLSL}
float16_t fmod(float16_t a, float16_t b);
float fmod(float a, float b);
template <uint N>
vector<float16_t, N> fmod(vector<float16_t, N> a, vector<float16_t, N> b);
template <uint N> vector<float, N> fmod(vector<float, N> a, vector<float, N> b);
\end{HLSL}

\begin{itemize}
  \item \textit{Returns:} A scalar or vector object of the same dimensions as
  the input where the value of each returned element is the floating-point
  remainder of the corresponding input elements.
\end{itemize}

\begin{HLSL}
float16_t frac(float16_t x);
float frac(float x);
template <uint N> vector<float16_t, N> frac(vector<float16_t, N> x);
template <uint N> vector<float, N> frac(vector<float, N> x);
\end{HLSL}

\begin{itemize}
  \item \textit{Returns:} A scalar or vector object of the same dimensions as
  the input where the value of each returned element is the floating-point
  fractional component of the corresponding input elements.
\end{itemize}

\begin{HLSL}
float frexp(float x, out float exp);
template <uint N>
vector<float, N> frexp(vector<float, N> x, out vector<float, N> exp);
\end{HLSL}

\begin{itemize}
  \item \textit{Returns:} A scalar or vector object of the same dimensions as
  the input where the value of each returned element is the mantissa of the
  corresponding input elements.
  \item \textit{Returns:} The scalar or vector \texttt{exp} is of the same
  dimensions as the input and the value of each element is initialized to the
  exponent of the corresponding input elements.
\end{itemize}

\begin{HLSL}
float16_t ldexp(float16_t x, float16_t exp);
float ldexp(float x, float exp);
template <uint N>
vector<float16_t, N> ldexp(vector<float16_t, N> x, vector<float16_t, N> exp);
template <uint N>
vector<float, N> ldexp(vector<float, N> x, vector<float, N> exp);
\end{HLSL}

\begin{itemize}
  \item \textit{Returns:} A scalar or vector object of the same dimensions as
  the input where the value of each returned element is the input element
  multiplied by the base 2 exponent of the corresponding element of \texttt{exp}
  (\(x_0 * 2^{exp_0}, .. x_n * 2^{exp_n}\)).
\end{itemize}

\begin{HLSL}
float16_t lerp(float16_t a, float16_t b, float16_t s);
float lerp(float a, float b, float s);
template <uint N>
vector<float16_t, N> lerp(vector<float16_t, N> a, vector<float16_t, N> b,
                          vector<float16_t, N> s);
template <uint N>
vector<float, N> lerp(vector<float, N> a, vector<float, N> b,
                      vector<float, N> s);
\end{HLSL}

\begin{itemize}
  \item \textit{Returns:} A scalar or vector object of the same dimensions as
  the input where the value of each returned element is the linear interpolation
  of the corresponding input \texttt{a} towards \texttt{b} by the step factor
  \texttt{s} (\(a_0 + s_0(b_0-a_0), .. a_0 + s_n(b_n-a_n)\)).
\end{itemize}

\begin{HLSL}
float16_t log(float16_t x);
float log(float x);
template <uint N> vector<float16_t, N> log(vector<float16_t, N> x);
template <uint N> vector<float, N> log(vector<float, N> x);
\end{HLSL}

\begin{itemize}
  \item \textit{Returns:} A scalar or vector object of the same dimensions as
  the input where the value of each returned element is the base $e$ logarithm
  of the corresponding element (\(\ln x_0, .. \ln x_n\)).
\end{itemize}

\begin{HLSL}
float16_t log10(float16_t x);
float log10(float x);
template <uint N> vector<float16_t, N> log10(vector<float16_t, N> x);
template <uint N> vector<float, N> log10(vector<float, N> x);
\end{HLSL}

\begin{itemize}
  \item \textit{Returns:} A scalar or vector object of the same dimensions as
  the input where the value of each returned element is the base 10 logarithm
  of the corresponding element (\(\log_{10} x_0, .. \log_{10} x_n\)).
\end{itemize}

\begin{HLSL}
float16_t log2(float16_t x);
float log2(float x);
template <uint N> vector<float16_t, N> log2(vector<float16_t, N> x);
template <uint N> vector<float, N> log2(vector<float, N> x);
\end{HLSL}

\begin{itemize}
  \item \textit{Returns:} A scalar or vector object of the same dimensions as
  the input where the value of each returned element is the base 2 logarithm
  of the corresponding element (\(\log_{2} x_0, .. \log_{2} x_n\)).
\end{itemize}

\begin{HLSL}
float16_t mad(float16_t a, float16_t b, float16_t c);
float mad(float a, float b, float c);
double mad(double a, double b, double c);
int16_t mad(int16_t a, int16_t b, int16_t c);
uint16_t mad(uint16_t a, uint16_t b, uint16_t c);
int mad(int a, int b, int c);
uint mad(uint a, uint b, uint c);
int64_t mad(int64_t a, int64_t b, int64_t c);
uint64_t mad(uint64_t a, uint64_t b, uint64_t c);
template <uint N>
vector<float16_t, N> mad(vector<float16_t, N> a, vector<float16_t, N> b,
                         vector<float16_t, N> c);
template <uint N>
vector<float, N> mad(vector<float, N> a, vector<float, N> b,
                     vector<float, N> c);
template <uint N>
vector<double, N> mad(vector<double, N> a, vector<double, N> b,
                      vector<double, N> c);
template <uint N>
vector<int16_t, N> mad(vector<int16_t, N> a, vector<int16_t, N> b,
                       vector<int16_t, N> c);
template <uint N>
vector<uint16_t, N> mad(vector<uint16_t, N> a, vector<uint16_t, N> b,
                        vector<uint16_t, N> c);
template <uint N>
vector<int, N> mad(vector<int, N> a, vector<int, N> b, vector<int, N> c);
template <uint N>
vector<uint, N> mad(vector<uint, N> a, vector<uint, N> b, vector<uint, N> c);
template <uint N>
vector<int64_t, N> mad(vector<int64_t, N> a, vector<int64_t, N> b,
                       vector<int64_t, N> c);
template <uint N>
vector<uint64_t, N> mad(vector<uint64_t, N> a, vector<uint64_t, N> b,
                        vector<uint64_t, N> c);
\end{HLSL}

\begin{itemize}
  \item \textit{Returns:} A scalar or vector object of the same dimensions as
  the input where the value of each returned element is a multiplication then
  addition of the corresponding input elements (\(a_0 * b_0 + c_0, .. a_n * b_n
  + c_n\)).
  \item \textit{Remarks:} An implementation may either fuse the multiplication
  and addition into a single operation resulting in a higher-precision result or
  it may implement the operation as separate multiplication and addition
  operations.
  \item \textit{Remarks:} In the presence of the \texttt{precise} keyword an
  implementation must always lower this operation to either a fused or a
  non-fused operation.
\end{itemize}

\begin{HLSL}
float16_t max(float16_t a, float16_t b);
float max(float a, float b);
double max(double a, double b);
int16_t max(int16_t a, int16_t b);
uint16_t max(uint16_t a, uint16_t b);
int max(int a, int b);
uint max(uint a, uint b);
int64_t max(int64_t a, int64_t b);
uint64_t max(uint64_t a, uint64_t b);
template <uint N>
vector<float16_t, N> max(vector<float16_t, N> a, vector<float16_t, N> b);
template <uint N> vector<float, N> max(vector<float, N> a, vector<float, N> b);
template <uint N>
vector<double, N> max(vector<double, N> a, vector<double, N> b);
template <uint N>
vector<int16_t, N> max(vector<int16_t, N> a, vector<int16_t, N> b);
template <uint N>
vector<uint16_t, N> max(vector<uint16_t, N> a, vector<uint16_t, N> b);
template <uint N> vector<int, N> max(vector<int, N> a, vector<int, N> b);
template <uint N> vector<uint, N> max(vector<uint, N> a, vector<uint, N> b);
template <uint N>
vector<int64_t, N> max(vector<int64_t, N> a, vector<int64_t, N> b);
template <uint N>
vector<uint64_t, N> max(vector<uint64_t, N> a, vector<uint64_t, N> b);
\end{HLSL}

\begin{itemize}
  \item \textit{Returns:} A scalar or vector object of the same dimensions as
  the input where the value of each returned element is the greater of the
  corresponding input elements (\(a_0 > b_0 ? a_0 : b_0, .. a_n > b_n ? a_n :
  b_n\)).
\end{itemize}

\begin{HLSL}
float16_t min(float16_t a, float16_t b);
float min(float a, float b);
double min(double a, double b);
int16_t min(int16_t a, int16_t b);
uint16_t min(uint16_t a, uint16_t b);
int min(int a, int b);
uint min(uint a, uint b);
int64_t min(int64_t a, int64_t b);
uint64_t min(uint64_t a, uint64_t b);
template <uint N>
vector<float16_t, N> min(vector<float16_t, N> a, vector<float16_t, N> b);
template <uint N> vector<float, N> min(vector<float, N> a, vector<float, N> b);
template <uint N>
vector<double, N> min(vector<double, N> a, vector<double, N> b);
template <uint N>
vector<int16_t, N> min(vector<int16_t, N> a, vector<int16_t, N> b);
template <uint N>
vector<uint16_t, N> min(vector<uint16_t, N> a, vector<uint16_t, N> b);
template <uint N> vector<int, N> min(vector<int, N> a, vector<int, N> b);
template <uint N> vector<uint, N> min(vector<uint, N> a, vector<uint, N> b);
template <uint N>
vector<int64_t, N> min(vector<int64_t, N> a, vector<int64_t, N> b);
template <uint N>
vector<uint64_t, N> min(vector<uint64_t, N> a, vector<uint64_t, N> b);
\end{HLSL}

\begin{itemize}
  \item \textit{Returns:} A scalar or vector object of the same dimensions as
  the input where the value of each returned element is the lesser of the
  corresponding input elements (\(a_0 < b_0 ? a_0 : b_0, .. a_n < b_n ? a_n :
  b_n\)).
\end{itemize}

\begin{HLSL}
float16_t pow(float16_t x, float16_t y);
float pow(float x, float y);
template <uint N>
vector<float16_t, N> pow(vector<float16_t, N> x, vector<float16_t, N> y);
template <uint N> vector<float, N> pow(vector<float, N> x, vector<float, N> y);
\end{HLSL}

\begin{itemize}
  \item \textit{Returns:} A scalar or vector object of the same dimensions as
  the input where the value of each returned element is the corresponding
  \texttt{x} element raised to the power of the corresponding \texttt{y} element
  (\(x_0^{y_0}, .. x_n^{y_n}\)).
\end{itemize}

\begin{HLSL}
float16_t rcp(float16_t x);
float rcp(float x);
double rcp(double x);
template <uint N> vector<float16_t, N> rcp(vector<float16_t, N> x);
template <uint N> vector<float, N> rcp(vector<float, N> x);
template <uint N> vector<double, N> rcp(vector<double, N> x);
\end{HLSL}

\begin{itemize}
  \item \textit{Returns:} A scalar or vector object of the same dimensions as
  the input where the value of each returned element is the reciprocal of the
  corresponding input element (\(1/x_0, .. 1/x_n\)).
\end{itemize}

\begin{HLSL}
float16_t round(float16_t x);
float round(float x);
template <uint N> vector<float16_t, N> round(vector<float16_t, N> x);
template <uint N> vector<float, N> round(vector<float, N> x);
\end{HLSL}

\begin{itemize}
  \item \textit{Returns:} A scalar or vector object of the same dimensions as
  the input where the value of each returned element is the value of the
  corresponding input element rounded to the nearest integer. Ties are rounded
  to the nearest even.
\end{itemize}

\begin{HLSL}
float16_t rsqrt(float16_t x);
float rsqrt(float x);
template <uint N> vector<float16_t, N> rsqrt(vector<float16_t, N> x);
template <uint N> vector<float, N> rsqrt(vector<float, N> x);
\end{HLSL}

\begin{itemize}
  \item \textit{Returns:} A scalar or vector object of the same dimensions as
  the input where the value of each returned element is the reciprocal of the
  square root of the corresponding input element (\(1/\sqrt{x_0}, .. 1/\sqrt{x_n}\)).
\end{itemize}

\begin{HLSL}
int sign(float16_t x);
int sign(float x);
int sign(double x);
int sign(int16_t x);
int sign(uint16_t x);
int sign(int x);
int sign(uint x);
int sign(int64_t x);
int sign(uint64_t x);
template <uint N> vector<int, N> sign(vector<float16_t, N> x);
template <uint N> vector<int, N> sign(vector<float, N> x);
template <uint N> vector<int, N> sign(vector<double, N> x);
template <uint N> vector<int, N> sign(vector<int16_t, N> x);
template <uint N> vector<int, N> sign(vector<uint16_t, N> x);
template <uint N> vector<int, N> sign(vector<int, N> x);
template <uint N> vector<int, N> sign(vector<uint, N> x);
template <uint N> vector<int, N> sign(vector<int64_t, N> x);
template <uint N> vector<int, N> sign(vector<uint64_t, N> x);
\end{HLSL}

\begin{itemize}
  \item \textit{Returns:} A scalar or vector object of the same dimensions as
  the input where the value of each returned element is the sign of the
  corresponding input elements. For an input element less than zero the result
  will be $-1$, for an input element greater than zero the result will be $1$,
  for an input element equal to zero the result will be $0$.
\end{itemize}

\begin{HLSL}
float16_t smoothstep(float16_t a, float16_t b, float16_t x);
float smoothstep(float a, float b, float x);
template <uint N>
vector<float16_t, N> smoothstep(vector<float16_t, N> a, vector<float16_t, N> b,
                                vector<float16_t, N> x);
template <uint N>
vector<float, N> smoothstep(vector<float, N> a, vector<float, N> b,
                            vector<float, N> x);
\end{HLSL}

\begin{itemize}
  \item \textit{Returns:} A scalar or vector object of the same dimensions as
  the input where the value of each returned element is the smooth Hermite
  interpolation of the corresponding input \texttt{a} towards \texttt{b} by the
  step factor
  \texttt{s} (\(a_0 + s_0(b_0-a_0), .. a_0 + s_n(b_n-a_n)\)).
\end{itemize}

\begin{HLSL}
float16_t sqrt(float16_t x);
float sqrt(float x);
template <uint N> vector<float16_t, N> sqrt(vector<float16_t, N> x);
template <uint N> vector<float, N> sqrt(vector<float, N> x);
\end{HLSL}

\begin{itemize}
  \item \textit{Returns:} A scalar or vector object of the same dimensions as
  the input where the value of each returned element is the square root of the
  corresponding input element (\(\sqrt{x_0}, .. \sqrt{x_n}\)).
\end{itemize}

\begin{HLSL}
float16_t step(float16_t a, float16_t x);
float step(float a, float x);
template <uint N>
vector<float16_t, N> step(vector<float16_t, N> a, vector<float16_t, N> x);
template <uint N> vector<float, N> step(vector<float, N> a, vector<float, N> x);
\end{HLSL}

\begin{itemize}
  \item \textit{Returns:} A scalar or vector object of the same dimensions as
  the input where the value of each returned element is the sign of the
  corresponding input elements. For an input element less than zero the result
  will be $-1$, for an input element greater than zero the result will be $1$,
  for an input element equal to zero the result will be $0$.

  \item \textit{Returns:} A scalar or vector object of the same dimensions as
  the input where the value of each returned element is $1$ if the \texttt{a}
  element is greater than or equal to thee \texttt{b} element; otherwise the
  element is $0$ (\(a_0 \geq b_0 ? 1 : 0, .. a_n \geq b_n ? 1 : 0\)).
\end{itemize}

\begin{HLSL}
float16_t trunc(float16_t x);
float trunc(float x);
template <uint N> vector<float16_t, N> trunc(vector<float16_t, N> x);
template <uint N> vector<float, N> trunc(vector<float, N> x);
\end{HLSL}

\begin{itemize}
  \item \textit{Returns:} A scalar or vector object of the same dimensions as
  the input where the value of each returned element is the truncated integer
  value of the corresponding input element. This rounds towards zero for all
  values, behaving like \texttt{floor} for positive values and \texttt{ceil} for
  negative values.
\end{itemize}

